\begin{abstract}
Tool-using language-model agents can turn untrusted text into externally committed actions. Existing defenses commonly classify prompts, plans, or whole tool calls, but one call can combine effects authorized for different resources, principals, modes, and provenance. We formulate \emph{tool-effect binding}: before commitment, a runtime must mediate a conservative causal abstraction of the security-relevant state transitions a call will realize. \sys{} separates offline counterfactual contract registration from deterministic online mediation. Registration intervenes on structured arguments, defaults, and trusted state, then rejects a contract when a real security-effect change is not represented by its atoms. Online, \sys{} totalizes defaults, instantiates the frozen atoms, and checks them against a task authority bound. A conditional confinement theorem exposes causal-abstraction soundness, envelope soundness, complete mediation, check--use integrity, and fail-closed uncertainty as explicit obligations. A field-level prototype completes all 726 official AgentDojo cases with 0/629 benchmark injection goals achieved, but only 0.330 benign and 0.304 attack-side task utility. An implementation audit exactly matches all 3,173 executed valid calls to pre-commit signatures. Compound-effect, authority-expansion, and omitted-default counterexamples keep the theorem conditional. Thus the current evidence supports exact sandbox mediation and a falsifiable effect-contract architecture, not production or unconditional trajectory safety.
\end{abstract}
